\section{Sections of quasi-coherent sheaves}
\label{section-sections-quasi-coherent}

\noindent
Here is a computation of sections of a quasi-coherent sheaf on a quasi-compact
open of an affine spectrum.

\begin{lemma}
\label{lemma-sections-over-quasi-compact-open-in-affine}
Let $A$ be a ring.
Let $I \subset A$ be a finitely generated ideal.
Let $M$ be an $A$-module.
Then there is a canonical map
$$
\colim_n \Hom_A(I^n, M)
\longrightarrow
\Gamma(\Spec(A) \setminus V(I), \widetilde{M}).
$$
This map is always injective.
If for all $x \in M$ we have $Ix = 0 \Rightarrow x = 0$
then this map is an isomorphism. In general, set
$M_n = \{x \in M \mid I^nx = 0\}$, then there is an
isomorphism
$$
\colim_n \Hom_A(I^n, M/M_n)
\longrightarrow
\Gamma(\Spec(A) \setminus V(I), \widetilde{M}).
$$
\end{lemma}

\begin{proof}
Since $I^{n + 1} \subset I^n$ and $M_n \subset M_{n + 1}$ we can
use composition via these maps to get canonical maps of $A$-modules
$$
\Hom_A(I^n, M)
\longrightarrow
\Hom_A(I^{n + 1}, M)
$$
and
$$
\Hom_A(I^n, M/M_n)
\longrightarrow
\Hom_A(I^{n + 1}, M/M_{n + 1})
$$
which we will use as the transition maps in the systems. Given an
$A$-module map $\varphi : I^n \to M$, then we get a map of
sheaves $\widetilde{\varphi} : \widetilde{I^n} \to \widetilde{M}$
which we can restrict to the open $\Spec(A) \setminus V(I)$.
Since $\widetilde{I^n}$ restricted to this open gives the structure
sheaf we get an element of
$\Gamma(\Spec(A) \setminus V(I), \widetilde{M})$.
We omit the verification that this is compatible with the transition maps
in the system $\Hom_A(I^n, M)$. This gives the first arrow.
To get the second arrow we note that
$\widetilde{M}$ and $\widetilde{M/M_n}$ agree over the open
$\Spec(A) \setminus V(I)$ since the sheaf $\widetilde{M_n}$
is clearly supported on $V(I)$. Hence we can use the same mechanism
as before.

\medskip\noindent
Next, we work out how to define this arrow in terms of algebra.
Say $I = (f_1, \ldots, f_t)$. Then
$\Spec(A) \setminus V(I) = \bigcup_{i = 1, \ldots, t} D(f_i)$.
Hence
$$
0 \to
\Gamma(\Spec(A) \setminus V(I), \widetilde{M}) \to
\bigoplus\nolimits_i M_{f_i} \to
\bigoplus\nolimits_{i, j} M_{f_if_j}
$$
is exact. Suppose that $\varphi : I^n \to M$ is an $A$-module map.
Consider the vector of elements $\varphi(f_i^n)/f_i^n \in M_{f_i}$.
It is easy to see that this vector maps to zero in the
second direct sum of the exact sequence above. Whence an element
of $\Gamma(\Spec(A) \setminus V(I), \widetilde{M})$.
We omit the verification that this description agrees with the one
given above.

\medskip\noindent
Let us show that the first arrow is injective using this description.
Namely, if $\varphi$ maps to zero, then for each $i$ the element
$\varphi(f_i^n)/f_i^n$ is zero in $M_{f_i}$. In other words we
see that for each $i$ we have $f_i^m\varphi(f_i^n) = 0$ for some $m \geq 0$.
We may choose a single $m$ which works for all $i$. Then we see that
$\varphi(f_i^{n + m}) = 0$ for all $i$. It is easy to see that
this means that $\varphi|_{I^{t(n + m - 1) + 1}} = 0$ in other
words that $\varphi$ maps to zero in the $t(n + m - 1) + 1$st
term of the colimit. Hence injectivity follows.

\medskip\noindent
Note that each $M_n = 0$ in case we have
$Ix = 0 \Rightarrow x = 0$ for $x \in M$. Thus
to finish the proof of the lemma it suffices to show that
the second arrow is an isomorphism.

\medskip\noindent
Let us attempt to construct an inverse of the second map of the lemma.
Let $s \in \Gamma(\Spec(A) \setminus V(I), \widetilde{M})$.
This corresponds to a vector $x_i/f_i^n$ with $x_i \in M$ of the
first direct sum of the exact sequence above.
Hence for each $i, j$ there exists $m \geq 0$
such that $f_i^m f_j^m (f_j^n x_i - f_i^n x_j) = 0$ in $M$.
We may choose a single $m$ which works for all pairs $i, j$.
After replacing $x_i$ by $f_i^mx_i$ and $n$ by $n + m$ we see
that we get $f_j^nx_i = f_i^nx_j$ in $M$ for all $i, j$.
Let us introduce
$$
K_n = \{x \in M \mid f_1^nx = \ldots = f_t^nx = 0\}
$$
We claim there is an $A$-module map
$$
\varphi :
I^{t(n - 1) + 1}
\longrightarrow
M/K_n
$$
which maps the monomial
$f_1^{e_1} \ldots f_t^{e_t}$ with $\sum e_i = t(n - 1) + 1$
to the class modulo $K_n$ of the expression
$f_1^{e_1} \ldots f_i^{e_i - n} \ldots f_t^{e_t}x_i$
where $i$ is chosen such that $e_i \geq n$ (note that there
is at least one such $i$).
To see that this is indeed the case suppose that
$$
\sum\nolimits_{E = (e_1, \ldots, e_t), |E| = t(n - 1) + 1}
a_E f_1^{e_1} \ldots f_t^{e_t} = 0
$$
is a relation between the monomials with coefficients $a_E$ in $A$.
Then we would map this to
$$
z =
\sum\nolimits_{E = (e_1, \ldots, e_t), |E| = t(n - 1) + 1}
a_E f_1^{e_1} \ldots f_{i(E)}^{e_{i(E)} - n} \ldots f_t^{e_t}x_{i(E)}
$$
where for each multiindex $E$ we have chosen a particular $i(E)$
such that $e_{i(E)} \geq n$.
Note that if we multiply this by $f_j^n$ for any $j$, then
we get zero, since by the relations $f_j^nx_i = f_i^nx_j$ above we get
\begin{align*}
f_j^nz & = \sum\nolimits_{E = (e_1, \ldots, e_t), |E| = t(n - 1) + 1}
a_E f_1^{e_1} \ldots f_j^{e_j + n}
\ldots f_{i(E)}^{e_{i(E)} - n} \ldots f_t^{e_t}x_{i(E)} \\
& =
\sum\nolimits_{E = (e_1, \ldots, e_t), |E| = t(n - 1) + 1}
a_E f_1^{e_1} \ldots f_t^{e_t}x_j
= 0.
\end{align*}
Hence $z \in K_n$ and we see that every relation gets mapped to zero
in $M/K_n$. This proves the claim.

\medskip\noindent
Note that $K_n \subset M_{t(n - 1) + 1}$. Hence the
map $\varphi$ in particular gives rise to an $A$-module map
$I^{t(n - 1) + 1} \to M/M_{t(n - 1) + 1}$.
This proves the second arrow of the lemma is surjective.
We omit the proof of injectivity.
\end{proof}

\begin{example}
\label{example-does-not-work-in-general}
We will give two examples showing that the first displayed map of
Lemma \ref{lemma-sections-over-quasi-compact-open-in-affine}
is not an isomorphism.

\medskip\noindent
Let $k$ be a field. Consider the ring
$$
A = k[x, y, z_1, z_2, \ldots]/(x^nz_n).
$$
Set $I = (x)$ and let $M = A$. Then the element $y/x$ defines a section of
the structure sheaf of $\Spec(A)$ over $D(x) = \Spec(A) \setminus V(I)$.
We claim that $y/x$ is not in the image of the canonical map
$\colim \Hom_A(I^n, A) \to A_x = \mathcal{O}(D(x))$. Namely, if so it
would come from a homomorphism $\varphi : I^n \to A$ for some $n$.
Set $a = \varphi(x^n)$. Then we would have $x^m(xa - x^ny) = 0$ for some
$m > 0$. This would mean that $x^{m + 1}a = x^{m + n}y$. This would mean that
$\varphi(x^{n + m + 1}) = x^{m + n}y$. This leads to a contradiction
because it would imply that
$$
0 = \varphi(0) = \varphi(z_{n + m + 1} x^{n + m + 1}) =
x^{m + n}y z_{n + m + 1}
$$
which is not true in the ring $A$.

\medskip\noindent
Let $k$ be a field. Consider the ring
$$
A = k[f, g, x, y, \{a_n, b_n\}_{n \geq 1}]/
(fy - gx, \{a_nf^n + b_ng^n\}_{n \geq 1}).
$$
Set $I = (f, g)$ and let $M = A$. Then $x/f \in A_f$ and $y/g \in A_g$
map to the same element of $A_{fg}$. Hence these define a section $s$
of the structure sheaf of $\Spec(A)$ over
$D(f) \cup D(g) = \Spec(A) \setminus V(I)$. However, there is no
$n \geq 0$ such that $s$ comes from an $A$-module map
$\varphi : I^n \to A$ as in the source of the first displayed arrow
of Lemma \ref{lemma-sections-over-quasi-compact-open-in-affine}.
Namely, given such a module map set
$x_n = \varphi(f^n)$ and $y_n = \varphi(g^n)$.
Then $f^mx_n = f^{n + m - 1}x$ and
$g^my_n = g^{n + m - 1}y$ for some $m \geq 0$ (see proof
of the lemma). But then we would have
$0 = \varphi(0) =
\varphi(a_{n + m}f^{n + m} + b_{n + m}g^{n + m}) =
a_{n + m}f^{n + m - 1}x + b_{n + m}g^{n + m - 1}y$ which is not the case
in the ring $A$.
\end{example}

\noindent
We will improve on the following lemma in the Noetherian case, see
Cohomology of Schemes, Lemma \ref{coherent-lemma-homs-over-open}.

\begin{lemma}
\label{lemma-sections-over-quasi-compact-open}
Let $X$ be a quasi-compact scheme.
Let $\mathcal{I} \subset \mathcal{O}_X$ be a
quasi-coherent sheaf of ideals of finite type.
Let $Z \subset X$ be the closed subscheme
defined by $\mathcal{I}$ and set $U = X \setminus Z$.
Let $\mathcal{F}$ be a quasi-coherent $\mathcal{O}_X$-module.
The canonical map
$$
\colim_n \Hom_{\mathcal{O}_X}(\mathcal{I}^n,
\mathcal{F})
\longrightarrow
\Gamma(U, \mathcal{F})
$$
is injective. Assume further that $X$ is quasi-separated.
Let $\mathcal{F}_n \subset \mathcal{F}$
be subsheaf of sections annihilated by $\mathcal{I}^n$.
The canonical map
$$
\colim_n \Hom_{\mathcal{O}_X}(\mathcal{I}^n,
\mathcal{F}/\mathcal{F}_n)
\longrightarrow
\Gamma(U, \mathcal{F})
$$
is an isomorphism.
\end{lemma}

\begin{proof}
Let $\Spec(A) = W \subset X$ be an affine open. Write
$\mathcal{F}|_W = \widetilde{M}$ for some $A$-module $M$
and $\mathcal{I}|_W = \widetilde{I}$ for some finite type
ideal $I \subset A$. Restricting the first displayed map
of the lemma to $W$ we obtain the first displayed map of
Lemma \ref{lemma-sections-over-quasi-compact-open-in-affine}.
Since we can cover $X$ by a finite number of affine opens this
proves the first displayed map of the lemma is injective.

\medskip\noindent
We have $\mathcal{F}_n|_W = \widetilde{M_n}$ where
$M_n \subset M$ is defined as in
Lemma \ref{lemma-sections-over-quasi-compact-open-in-affine}
(details omitted). The lemma guarantees that we have a bijection
$$
\colim_n  \Hom_{\mathcal{O}_W}(
\mathcal{I}^n|_W, (\mathcal{F}/\mathcal{F}_n)|_W)
\longrightarrow
\Gamma(U \cap W, \mathcal{F})
$$
for any such affine open $W$.

\medskip\noindent
To see the second displayed arrow of the lemma is bijective,
we choose a finite affine open covering $X = \bigcup_{j = 1, \ldots, m} W_j$.
The injectivity follows immediately from the above and the finiteness of
the covering. If $X$ is quasi-separated, then for each pair
$j, j'$ we choose a finite affine open covering
$$
W_j \cap W_{j'} = \bigcup\nolimits_{k = 1, \ldots, m_{jj'}} W_{jj'k}.
$$
Let $s \in \Gamma(U, \mathcal{F})$. As seen above for each $j$ there exists
an $n_j$ and a map
$\varphi_j : \mathcal{I}^{n_j}|_{W_j} \to
(\mathcal{F}/\mathcal{F}_{n_j})|_{W_j}$
which corresponds to $s|_{U \cap W_j}$.
By the same token for each triple $(j, j', k)$ there exists an integer
$n_{jj'k}$ such that the restriction of $\varphi_j$ and $\varphi_{j'}$
as maps $\mathcal{I}^{n_{jj'k}} \to \mathcal{F}/\mathcal{F}_{n_{jj'k}}$
agree over $W_{jj'k}$. Let $n = \max\{n_j, n_{jj'k}\}$ and we see that
the $\varphi_j$ glue as maps
$\mathcal{I}^n \to \mathcal{F}/\mathcal{F}_n$ over $X$.
This proves surjectivity of the map.
\end{proof}













